\begin{enumerate}
    \item \emp{Стиль изложения.}
        Каждое слово и предложение должно нести конкретную математическую или учебно-методическую нагрузку (<<если удалить это слово или предложение, потеряет ли читатель что-то существенное?>>).
        Незначимые слова и предложения удаляются без сожалений.

    \item \emp{Собственные команды.}
        Если в файл \verb|helpers.tex| добавляется новая команда, она сразу же должна быть описана в подходящем разделе данного руководства по стилю оформления.

    \item \emp{Заметки.}
        Для задания примечаний на этапе разработки книги используются команды \verb|\td{текст}| (от англ. <<todo>>; рекомендации по добавлению текста) и \verb|\fx{текст}| (от англ. <<fix>>; рекомендации по исправлению текста).

    \item \emp{Предложения.}
        В исходном коде делается перенос строки после каждого предложения, так как это упрощает отслеживание изменений в системе контроля версий и облегчает выявление ошибок при сборке документа.
        Внутри одного предложения дополнительные переводы строки по возможности не используются (но допустимы, если это в существенной мере улучшает читабельность предложения).

    \item \emp{Дефисы и тире.}
        Используются три различных знака (замена одного знака другим не допускается), согласно нормам типографики:
        \begin{itemize}
            \item Дефис <<->>
                (набор: \verb|-|)~---
                внутри слов и при дефисном написании, без пробелов: <<физико-математический>>, <<всё-таки>>, <<кое-что>>.
            \item Короткое тире <<-->>
                (набор: \verb|--|)~---
                для числовых и именованных диапазонов, без пробелов: <<1--10>>, <<теорема Больцано--Вейерштрасса>>.
                Используется также при указании страниц (<<\verb|с.~12--15|>>) и интервалов времени.
            \item Длинное тире <<--->>
                (набор: \verb|---|)~---
                пунктуационное тире в предложении, с обычным пробелом справа и неразрывным пробелом слева (чтобы тире не начинало строку при переносе): <<\verb|Математика~--- царица наук|>>.
        \end{itemize}

    \item \emp{Многоточия.}
        Используется команда \verb|\ldots| вместо вручную набранных трёх точек (например, $1, 2, \ldots, 42$).
        Если нужно центрированное выравнивание (например, $\sum_{i_1} \sum_{i_2} \cdots \sum_{i_d}$), то используется команда \verb|\cdots|.

    \item \emp{Кавычки.}
        Используются исключительно <<ёлочки>> (в коде: \verb|<<ёлочки>>|).
        Если возникает необходимость во вложенных кавычках, то делается так: \verb|<<снаружи ``вложение'' снаружи>>|.

    \item \emp{Буква ё.}
        Данная буква (не буква <<е>>) пишется везде, где должна быть по правилам русского языка.

    \item \emp{Сноски.}
        Сноски ставятся после знака препинания или слова, к которому относятся, без пробела (например: <<\verb|видно,\footnote{...}| что ...>>).
        Содержание сноски должно представлять собой полноценное предложение с соответствующими знаками препинания.

    \item \emp{Метки.}
        Задаётся префикс с двоеточием, указывающий на тип объекта:
        \verb|defi:name|~--- определения;
        \verb|stat:name|~---  утверждения;
        \verb|lemm:name|~---  леммы;
        \verb|theo:name|~---  теоремы;
        \verb|conc:name|~---  следствия;
        \verb|exam:name|~---  примеры;
        \verb|note:name|~---  замечания;
        \verb|algo:name|~---  алгоритмы;
        \verb|code:name|~---  листинги;
        \verb|sec:name|~---  разделы;
        \verb|fig:name|~---  рисунки;
        \verb|tbl:name|~---  таблицы;
        \verb|eq:name|~---  формулы.
        Используется как английский, так и русский язык для имён в метках, при этом пробелы в именах не допускаются и заменяются на нижние подчеркивания.
        Например: <<\verb|как следует из теоремы~\ref{theo:теорема_пифагора}|\ldots>>.

    \item \emp{Ссылки.}
        Для ссылки на формулу используется команда \verb|\eqref|, в остальных случаях~--- команда \verb|\ref|.
        Всегда ставится тильда (неразрывный пробел) перед ссылкой (например, \verb|см.~рис.~\ref{fig:name}| или \verb|см.~формулу~\eqref{eq:name}|); это запрещает перенос номера на следующую строку.

    \item \emp{Неразрывные пробелы} (\verb|~|).
        Используются в следующих случаях:
        \begin{itemize}
            \item перед командой ссылки: \verb|рис.~\ref{...}|, \verb|теоремы~\ref{...}|,
                \verb|формулы~\eqref{...}|;
            \item между сокращением и числом или именем: \verb|рис.~5|, \verb|стр.~7|;
            \item в сокращённых оборотах: \verb|и~т.п.|, \verb|и~т.д.|,
                \verb|и~др.|, \verb|и~пр.|;
            \item между инициалами и фамилией: \verb|А.~Н.~Колмогоров|;
            \item между числом и единицей измерения: \verb|5~м|, \verb|42~с|, \verb|10~кг|;
            \item перед знаком процента: \verb|5~\%|;
            \item в датах: \verb|1~января 2026~года|.
        \end{itemize}

    \item \emp{Числа в тексте.}
        Обычные числа в тексте набираются без математического режима.
        Например, записывается <<1, 2, \ldots, 10 глав>>, а не <<\$1\$, \$2\$, \ldots, \$10\$ глав>>.

    \item \emp{Пунктуация в пунктах списка.}
        Оформление пунктов согласуется с их синтаксической природой.
        Если пункт представляет полное предложение, то он начинается с заглавной
                буквы и заканчивается точкой.
        Если пункт является продолжением фразы, введённой обобщающим словом перед списком, то он начинается со строчной буквы и заканчивается точкой с запятой, а последний пункт заканчивается точкой.
        Если пункт содержит одно слово или короткое словосочетание, то допустима запятая как разделитель, при этом последний пункт заканчивается точкой.
        В рамках одного списка все пункты должны следовать одной модели.

    \item \emp{Окружения для формул.}
        Для внутритекстовой формулы используется стандартное окружение \verb|$...$|.
        Выключные формулы (то есть формулы, выносимые на отдельную строку), имеющие номер, оформляются через окружение \verb|equation|.
        Выключные формулы без номера оформляются через окружение \verb|equation*|.
        Прочие display-окружения верхнего уровня (\verb|align|, \verb|gather|, \verb|multline|, а также конструкции вида \verb|$$...$$|) не используются.

    \item \emp{Пунктуация в формулах.}
        Формула является частью предложения и подчиняется соответствующим правилам пунктуации:
        \begin{itemize}
            \item Двоеточие перед формулой ставится, если перед ней стоит обобщающее слово или вводящий оборот (<<имеем следующее:>>, <<определяется так:>>, <<справедливо:>>).
                Если же формула является продолжением фразы, то двоеточие не ставится.
            \item Знак препинания после выключной формулы ставится сразу за последним символом формулы (без пробела) с учётом правил русского языка: точка~--- если формула завершает предложение, запятая~--- если за ней следует придаточное предложение или продолжение перечисления, точка с запятой~--- если формула завершает пункт перечисления.
                Для внутритекстовой формулы знак препинания ставится после закрывающего символа \verb|$|.
            \item Если текст после внутритекстовой или выключной формулы является грамматическим продолжением предложения и не требует запятой или иного знака по правилам русского языка, то знак препинания не ставится.
                Примеры: <<Из равенства \ldots\ следует>>, <<Рассмотрим выражение \ldots\ и покажем>>, <<Действительно: \ldots\ ввиду того, что>>.
        \end{itemize}

    \item \emp{Пробелы в формулах.}
        Для ручной регулировки горизонтальных интервалов в математическом режиме используются две команды:
        \begin{itemize}
            \item \verb|\;|~--- пробел в 5/18~em; применяется для небольших отступов внутри формулы, например, для визуального разделения однотипных выражений в одной строке ($a < b, \; b < c$);
            \item \verb|\quad|~--- пробел в 1~em; применяется только в выключных формулах для разнесения крупных частей: вокруг словесных вставок ($x > 0 \quad \text{и} \quad y > 0$), между независимыми утверждениями в одной строке или перед пояснением справа.
        \end{itemize}
        Прочие команды пробелов (<<\verb|\,|>>, <<\verb|\:|>>, <<\verb|\ |>>, <<\verb|\qquad|>>, <<\verb|\!|>>) в основном тексте книги не применяются (но могут использоваться в \verb|helpers.tex| при определении новых команд).
        
    \item \emp{Текст в формулах.}
        Команда \verb|\mathrm{...}| используется для латинских многобуквенных математических идентификаторов и констант, которые должны набираться прямым шрифтом (например, $x_{\mathrm{max}}$, $\mathrm{const}$, $\mathrm{e}$, $\mathrm{i}$ для мнимой единицы), но не для вставки словесных комментариев, так как у \verb|\mathrm| нет правильных пробелов и переносов.
        Для вставки словесных фрагментов (как русских, так и любых, требующих корректных пробелов и переносов) в математический режим используется команда \verb|\text{...}|, которая сохраняет шрифт и интервалы окружающего абзаца.
        Например: $x \in A, \text{ где } A \subset \R$ (\verb|x \in A, \text{ где } A \subset \R|).
        Пробелы (обычные) вокруг \verb|\text{...}| должны входить внутрь аргумента, а не ставиться снаружи (если требуются пробелы вида \verb|\;| и \verb|\quad|, то они указываются вне блока \verb|\text{...}|).

    \item \emp{Нумерация формул.}
        Нумеруются только те формулы, на которые имеются ссылки в тексте.
        Для таких формул используется \verb|equation| с меткой \verb|\label{eq:name}| (метка ставится сразу после \verb|\begin{equation}| без пробела и на той же строке).
        Во всех остальных случаях применяется \verb|equation*| (без номера).

    \item \emp{Абзацы после выключных формул.}
        Не добавляется пустая строка после \verb|\end{equation}| (или \verb|\end{equation*}|), если следующий за формулой текст грамматически продолжает предложение или представляет следующее предложение того же самого абзаца (иначе LaTeX начнёт новый абзац с красной строки).
        Пустая строка добавляется только если после формулы начинается новый абзац.

    \item \emp{Знак завершения доказательства.}
        В блоках с доказательством или решением знак завершения доказательства добавляется автоматически.
        Поэтому вручную команды \verb|\qed|, \verb|\qedsymbol|, \verb|\hfill\qedsymbol| и аналогичные конструкции не ставятся.
        Если доказательство или решение заканчивается обычным текстом, дополнительных действий не требуется.
        Однако, если последним содержательным элементом доказательства или решения является выключная формула, то для размещения знака завершения доказательства на строке этой формулы следует использовать команду \verb|\qedhere| после завершающего знака препинания.

    \item \emp{Переносы в формулах.}
        Используется окружение \verb|split| внутри \verb|equation|.
        Знак бинарной операции или отношения (\verb|=|, \verb|+|,
        \verb|\le|, \verb|\im|), по которому происходит перенос, повторяется в начале новой строки.
        Если в конце строки перед \verb|\\| остаётся знак бинарной операции (\verb|+|, \verb|-|, \verb|\pm|, \verb|\mul| и~т.п.), то после него ставится пустая группа \verb|{}|: иначе TeX, не находя справа операнда, уберёт отступ, и знак <<прилипнет>> к предыдущему символу.
        Для знаков отношения (\verb|=|, \verb|\le|, \verb|\ge|, \verb|\im|, \verb|\eq| и~т.п.) такая правка не требуется.
        Например:
        \begin{equation*}
            \begin{split}
                (a + b)^2 
                    & = (a+b)(a+b) = \\
                    & = a^2 + ab + {} \\
                    & + ba + b^2 = a^2 + 2ab + b^2.
            \end{split}
        \end{equation*}

    \item \emp{Системы и группы формул.}
        Для одновременного представления нескольких связанных формул (система уравнений, цепочка преобразований, группа тождеств, выкладка с пояснениями справа) используется окружение \verb|aligned| внутри \verb|equation| или \verb|equation*|.
        В отличие от \verb|split|, окружение \verb|aligned| допускает несколько точек выравнивания: символы \verb|&| ставятся парами (первый \verb|&| в строке~--- точка выравнивания, второй~--- разделитель колонок).
        Пример цепочки преобразований с пояснениями справа:
        \begin{equation*}
            \begin{aligned}
                \lim_{n \to \infty} \frac{n^2 + 1}{2n^2 - n}
                    & = \lim_{n \to \infty} \frac{1 + 1/n^2}{2 - 1/n}
                        && \text{(делим на $n^2$)} \\
                    & = \frac{1 + 0}{2 - 0} = \frac{1}{2}.
                        && \text{(так как $1/n^2 \to 0$ и $1/n \to 0$)}
            \end{aligned}
        \end{equation*}

    \item \emp{Кусочные определения.}
        Кусочно-заданные функции и выражения оформляются окружением \verb|cases| внутри \verb|equation| (или \verb|equation*|).
        После каждой ветви через \verb|&| указывается условие, ветви разделяются \verb|\\|; знаки препинания внутри ветвей подчиняются общим правилам пунктуации формул.
        Например:
        \begin{equation*}
            \sgn(x) = \begin{cases}
                1,   & x > 0, \\
                 0,  & x = 0, \\
                -1,  & x < 0.
            \end{cases}
        \end{equation*}

    \item \emp{Рисунки и таблицы.}
        Оформляются окружениями \verb|figure| и \verb|table| соответственно.
        При этом:
        \begin{itemize}
            \item Плавающему окружению задаётся спецификатор позиционирования \verb|[t!]| (объект размещается в верхней части страницы; восклицательный знак ослабляет внутренние ограничения LaTeX на размещение плавающих объектов).
                Прочие спецификаторы (\verb|[h]|, \verb|[b]|, \verb|[p]|, \verb|[H]|) применяются в исключительных случаях.
            \item Подпись (полное предложение) размещается под рисунком (команда \verb|\caption| после содержимого) и над таблицей (до \verb|\begin{tabular}|).
                Пример: <<График функции $y = x^2$.>>
            \item Метка вида \verb|\label{fig:name}| или \verb|\label{tbl:name}| ставится после \verb|\caption| на следующей строке.
        \end{itemize}

    \item \emp{Определения величин и предикатов.}
        Для ввода новых понятий строго выполняется семантическое разделение между присваиванием значений и равносильностью утверждений:
        \begin{itemize}
            \item Для ввода новых алгебраических величин (чисел, переменных, функций, множеств) используется знак $=$ с двоеточием (равенство по определению).
                Применяется команда \verb|\is|, если определяемая величина находится слева ($A_{\mathrm{new}} \is B$) и команда \verb|\isr|, если определяемая величина находится справа ($C \isr D_{\mathrm{new}}$).
            \item Для ввода новых логических утверждений, свойств или предикатов используется знак $\Leftrightarrow$ с двоеточием (равносильность по определению).
                Применяются команды \verb|\eqdef| ($P(x) \eqdef (x > 0)$) и \verb|\eqdefr| ($(x \in A \land x \in B) \eqdefr Q(x)$).
        \end{itemize}
        Если это не ведёт к неоднозначности, то для задания промежуточных локальных величин в выкладках допустимо (но не желательно) использовать обычный знак $=$ или $\Leftrightarrow$.

    \item \emp{Локальная регулировка высоты страницы.}
        Используется стиль \verb|\raggedbottom|, при котором низ страницы не выравнивается.
        В отдельных случаях вёрстка разворота может оказаться визуально неудачной, тогда допустима ручная локальная правка с помощью команды \verb|\enlargethispage{\baselineskip}|, которая увеличивает высоту текущей страницы ровно на одну строку.
        Аналогично, \verb|\enlargethispage{-\baselineskip}| укорачивает текущую страницу на одну строку.
        Команда \verb|\enlargethispage| ставится в любом месте текущей страницы, её действие не распространяется на другие страницы и не меняет глобальные настройки геометрии.

\end{enumerate}